\bookchapter{Kapitel 0: Hallo, bevor es losgeht}

Bevor wir richtig loslegen, klären wir ein paar Grundregeln, damit du alles Spätere verstehst. Ganz ohne Eile: Wer das Werkzeug schärft, hackt danach schneller.

\booksection{Was dieses Buch für dich tun kann (und was nicht)}

Es hilft dir, systematisch zu verstehen, wie dein Smartphone funktioniert und welche Logik hinter den Funktionen steckt. Ob Android- oder Apple-Lager: Die meisten Denkmuster sind universell. Und es begleitet dich beim Ausprobieren – vom „Bedienen“ zum „Verstehen“.

Es ist kein Handbuch für ein bestimmtes Modell, kein Tech-News-Ticker und schon gar kein Kurs fürs Videos-Würgen. Ehrlich gesagt: Wer „noch komfortabler durch Kurzvideos scrollen“ lernen will, wird enttäuscht – uns geht es darum, dass das Handy wirklich dir dient, nicht umgekehrt.

\booksection{Unsere „Geheimzeichen“-Erklärung}

Damit alles eindeutig bleibt, benutzt das Buch einige kleine Markierungen – wie verabredete Geheimzeichen zwischen uns:

\begin{itemize}
  \item \textbf{[Schaltflächenname]}: Steht für konkrete Schaltflächen oder App-Namen auf dem Bildschirm, z. B. [Einstellungen], [OK].
  \item \textbf{>}: Zeigt die Reihenfolge der Schritte an – „tippe auf [Einstellungen] > [Ton]“ heißt: erst Einstellungen, dann Ton.
  \item \textbf{Fettdruck}: Sehr wichtig! Steht für zentrale Konzepte oder Sicherheitswarnungen – unbedingt beachten.
  \item \emph{Kursiv}: kleine ergänzende Gedanken oder Vergleiche, die dir das Verständnis erleichtern.
\end{itemize}

\booksection{Android und iOS, beide dabei}

Rate mal: Manche Dinge heißen anders, tun aber dasselbe. Z. B. Anwendung (Android) = App (iOS), Einstellungen (beide gleich), Dateimanager (Android) und Dateien (iOS). Hauptsache, du weißt, es ist dasselbe – lass dich von den Namen nicht verwirren.

Außerdem gibt es im Android-Lager noch Huaweis HarmonyOS, das eigene Besonderheiten hat – darauf gehen wir später gesondert ein. Kurzum: egal welches Smartphone du nutzt, dieses Buch hilft dir weiter.

\booksection{Handyspeicher: warum von 128 GB nach dem Kauf kaum 110 GB übrig bleiben}

Die Einheiten kennst du: 1 TB = 1024 GB, 1 GB = 1024 MB, 1 MB = 1024 KB … Aber hier lauert die Falle: Die Hersteller rechnen mit dem Faktor 1000, das System mit 1024 – deshalb „schrumpft“ der gekaufte 128-GB-Speicher auf dem Papier.

Dass der Speicher deines Handys schrumpft, hat vor allem drei Gründe:

\begin{enumerate}
  \item \textbf{Das System selbst belegt Platz}: Das Betriebssystem ist wie die tragenden Wände eines Hauses – es muss da sein. Aktuelle Systeme belegen oft 10-20 GB, das ist völlig normal.
  \item \textbf{Vom Hersteller vorinstallierte Apps}: Die vorinstallierten Apps auf einem neuen Smartphone belegen ebenfalls Platz – manche lassen sich deinstallieren, manche nicht.
  \item \textbf{Partitionierungs- und Formatierungsverlust}: Technische Gründe, so wie das Futter eines Rucksacks auch Platz braucht.
\end{enumerate}

\begin{tcbinfo}{Kleiner Tipp für den Smartphone-Kauf}
128 GB sind heute kaum noch ausreichend, 256 GB sind ein guter Start. Wer gerne fotografiert oder spielt, sollte direkt zu 512 GB oder mehr greifen. Denk dran: Lieber etwas mehr Speicher als zu wenig.

\end{tcbinfo}

\begin{tcbpractice}{Mach mit: Schau dir den "Ausweis" deines Smartphones an}
Öffne [Einstellungen], suche “Über das Telefon” (oder ähnlich) und notiere: Modell, Systemversion, Speicherplatz und freier Speicher. Das ist der erste Schritt, um es kennenzulernen. Wirf auch einen Blick auf die Akkugesundheit, damit du Bescheid weißt.

\end{tcbpractice}

\booksection{Dein Smartphone ist eine “digitale Vermögens”-Bibliothek}

Bevor wir ins Detail gehen, wollen wir eine Kern-Erkenntnis festhalten:\textbf{Das Smartphone ist längst mehr als ein Kommunikationsgerät – es ist unsere wichtigste digitale Vermögensbibliothek}. Kontakte, Chatverläufe, Familienfotos, Arbeitsdokumente, Banking-Apps, QR-Zahlcodes, Gesundheitsdaten, digitale Ausweise … zusammengenommen ist dieser Datenwert oft höher als der des Geräts selbst.

\begin{tcbinfo}{Warum ist das wichtig?}
Wenn das Smartphone verloren geht oder kaputt ist, verlierst du nicht nur Hardware im Wert von einigen tausend Yuan.\textbf{Oft sind es über Jahre angesammelte, nicht erneuerbare Daten}– Wachstumsfotos der Kinder, Sprachnachrichten verstorbener Angehöriger, wichtige Verträge ohne Backup. Strukturierte Pflege- und Schutzgewohnheiten schützen dich auf drei Ebenen:

\begin{itemize}
  \item \textbf{Höhere Effizienz}: Eine saubere Dateistruktur beschleunigt die Suche – das Smartphone bleibt langfristig flüssig.
  \item \textbf{Längere Lebensdauer}: Vernünftiges Speichermanagement entlastet Lese-/Schreibvorgänge, Akku und Speicher halten länger.
  \item \textbf{Datensicherheit}: Ebenen-Verwaltung + regelmäßige Backups + minimierte Berechtigungen halten "Datenkatastrophen" fern.
\end{itemize}

\end{tcbinfo}

\booksubsection{Unsere Kernphilosophie: “Prävention zuerst, Behandlung kombiniert”}

Wie bei der Gesundheit man nicht erst zum Arzt geht, wenn man krank ist.\textbf{Auch das Smartphone-Management beginnt bei “täglichen Gewohnheiten”}. Dieses Buch betont nicht irgendwelche Profi-Tricks, sondern diese schlichten Prinzipien:

\begin{itemize}
  \item \textbf{Prävention zuerst}: Apps nur aus offiziellen Stores, Fotos vor dem Teilen anonymisieren, keine doppelten Passwörter – diese "kleinen Handgriffe" kosten kaum Zeit, wehren aber die meisten Risiken ab.
  \item \textbf{Regelmäßige Kontrollen}: Einmal pro Monat Download-Ordner aufräumen, einmal pro Quartal Berechtigungen prüfen, alle sechs Monate ein Voll-Backup – ein fester Rhythmus ist hundertmal wirksamer als Notfall-Fehlerbehebung.
  \item \textbf{Puffer einplanen}: Cloud-Sync aktivieren, Zugangsdaten sicher notieren, eine "digitale Erbe-Liste" anlegen – damit du und deine Familie bei einem Notfall ruhig reagieren könnt.
\end{itemize}
